\setcounter{section}{0}
\section{自然数と加算集合}
\label{section:自然数と加算集合}

\begin{definition}[p16]
	\label{definition:集合論的な自然数の定義}
	以下のように,自然数$n$を集合論的に構成できる.\\
	$n = n - 1 \cup \{n - 1\}$と定義する.すると,
	$n = \{0,1,2,... n-1\}と表せる.$
\end{definition}

\begin{definition}[p183の定義7.1.1]
	\label{definition:有限集合,無限集合の定義}
	$X$を集合とする。$X$が有限集合,あるいは無限集合であることを以下のように定義する.\\
	 $X$は有限集合$\defarr \exists n \in \mathbb{N}\exists f( \; \BijectionMap{f}{n}{X} \; )$\\
	 $X$は無限集合$\defarr X$が無限集合でない.\\
\end{definition}

\begin{proposition}[p184の命題7.1.2]
	1. $n,m$を自然数とし,$f:n \rightarrow m$を全単射とする.このとき,
	$n=m$である.論理式は以下.\\
	$\forall n ,m \in \mathbb{N} (\exists f(\BijectionMap{f}{n}{m}) \rightarrow n = m)$\\
	2. 自然数全体の集合$\mathbb{N}$は無限集合である.論理式は以下.\\
	$\forall n \in \mathbb{N} \forall f( \lnot \; \BijectionMap{f}{n}{\mathbb{N}} \; )$
\end{proposition}

\begin{proof} 1. $\varphi (n) :\Leftrightarrow \forall m \in \mathbb{N} (\exists f(\BijectionMap{f}{n}{m}) \rightarrow n = m)$とし、$n$に関する数学的帰納法で示す.\\

(goal) $\varphi (0) \land \forall n \in \mathbb{N} ( \varphi (n) \rightarrow \varphi (n+1))$\\

(base step) $\varphi (0) \Leftrightarrow  \forall m \in \mathbb{N}  (\exists f(\BijectionMap{f}{0}{m}) \rightarrow 0 = m)$ \\
任意に$m \in \mathbb{N}$をとる.また,$\exists f(\BijectionMap{f}{0}{m})$は真より,そんな$f$を一つとり固定.\\
$m \neq 0$とし,背理法で示す.\\
$ 0 = \emptyset$より,$\BijectionMap{f}{\emptyset}{m}$は真.よって,$\SurjectionMap{f}{\emptyset}{m}$も真なので,以下が成り立つ.\\
$\forall y \in m \exists x \in \emptyset (f(x) = y)$\\
今、$m \neq \emptyset$なので,$\exists x \in \emptyset (f(x) = y)$が真であるが,
これは,$\forall x \in \emptyset$が偽であることに矛盾する.よって,goalAは真.\\

(induction step) $\forall n \in \mathbb{N} ( \varphi (n) \rightarrow \varphi (n+1)) \\
\Leftrightarrow \forall n \in \mathbb{N}(\forall m' \in \mathbb{N}  (\exists f'(\BijectionMap{f'}{n}{m'}) \rightarrow  n = m') \rightarrow 
\forall m \in \mathbb{N}  (\exists f(\BijectionMap{f}{n+1}{m}) \rightarrow  n+1 = m))$ \\
任意に$n,m \in \mathbb{N}$をとる.\\
また,$\exists f(\BijectionMap{f}{n+1}{m})$は真より,そんな$f$を一つとり固定.\\
このとき、$m \neq 0$,つまり$m-1 \in \mathbb{N}$である.\\

\begin{sub-block}{{\boldmath $\because$}\,}
\begin{itemize}
\item[] \hspace{-0.5cm} \\
	いま、$\BijectionMap{f}{n+1}{m}$が真で、$n \in \mathbb{N}$より、$n+1 \neq 0$.\\
　$m = 0$とすると、$\BijectionMap{f}{n+1}{m}$が成り立たないので、$m \neq 0$
\end{itemize}
\end{sub-block}
$m-1 \in \mathbb{N}$と,帰納法の仮定より,$ \exists f'(\BijectionMap{f'}{n}{m-1})\rightarrow  n = m-1$が真.\\

$l:=f(n)$とおく.\\
$g:=\{<k,k> \mid k \neq l, m-1 \} \cup \{<k,l> \mid k = m-1 \} \cup \{<k,m-1> \mid k = l \}$とする.\\
$g \circ g = id_m$なので,$\BijectionMap{g \circ g}{m}{m}$は真.よって,$\BijectionMap{g}{m}{m}$も真.\\

\begin{sub-block}{{\boldmath $\because$}\,}
\begin{itemize}
\item[] \hspace{-0.5cm} \\
	$\BijectionMap{g \circ g}{m}{m}$は真より,$\SurjectionMap{g \circ g}{m}{m} \land \InjectionMap{g \circ g}{m}{m}$.\\
	このとき,$\SurjectionMap{g}{m}{m} \land \InjectionMap{g}{m}{m}$は真より,(問題集)$\BijectionMap{g}{m}{m}$も真.
\end{itemize}
\end{sub-block}

また,$\BijectionMap{f}{n+1}{m}, \BijectionMap{g}{m}{m}$より,$\BijectionMap{g \circ f}{n+1}{m}$も真.\\
$f':= g \circ f \setminus \{<n, m-1>\}$とおく.\\
$g \circ f(n) = g(l) = m-1$より,$\BijectionMap{f'}{n}{m-1}$は真.\\
$f'$の存在より,$\exists f'(\BijectionMap{f'}{n}{m-1})$は真.ゆえに,$n = m-1$は真.\\
$n+1=n \cup \{n\}, m=m-1 \cup \{m-1\}$と表せる.\\
これと,$n=m-1$より,$n+1=m$.よってgoalBは真.\\
goalA $\land$ goalBが真より、goalは真.\\
\end{proof}

\begin{proof} 2. $\varphi (n) :\Leftrightarrow \forall f( \lnot \; \BijectionMap{f}{n}{\mathbb{N}} \; )$とし、$n$に関する数学的帰納法で示す.\\
(goal) $\varphi (0) \land \forall n \in \mathbb{N} ( \varphi (n) \rightarrow \varphi (n+1))$\\

(base step) $\varphi (0) \Leftrightarrow  \forall f( \lnot \; \BijectionMap{f}{\emptyset}{\mathbb{N}} \; )$ \\
背理法で示す.よって、$\exists f(  \; \BijectionMap{f}{\emptyset}{\mathbb{N}} \; )$が真とする.そんな$f$を一つとり固定する.\\
$\SurjectionMap{f}{\emptyset}{\mathbb{N}}$も真なので,以下が成り立つ.\\
$\forall y \in \mathbb{N} \exists x \in \emptyset (f(x) = y)$\\
$\mathbb{N} \neq \emptyset$なので,
$\exists y'(y' \in \mathbb{N})$は真.そんな$y'$を1つとり固定する.\\
$\exists x \in \emptyset (f(x) = y')$が真であるが,
これは,$\exists x (x \not\in \emptyset)$が真であることに矛盾する.\\

(induction step) $\forall n \in \mathbb{N} ( \varphi (n) \rightarrow \varphi (n+1)) 
\Leftrightarrow \forall n \in \mathbb{N} ( \forall f_n( \lnot \; \BijectionMap{f_n}{n}{\mathbb{N}} \; ) \rightarrow \forall f( \lnot \; \BijectionMap{f_{n+1}}{n+1}{\mathbb{N}} \; ))$\\
任意に$n \in \mathbb{N}$をとる.\\
任意に$f_n(\lnot \BijectionMap{f_n}{n}{\mathbb{N}})$をとる.\\
背理法で示すので,$\exists f_{n+1}(\; \BijectionMap{f_{n+1}}{n+1}{\mathbb{N}} \; )$が真とする.そんな$f$を一つとり固定.\\
$m := f_{n+1}(n)$とおく.\\
$g:=\{<k,k> \mid k \leq m-1 \land k \in \mathbb{N}\} \cup \{<k,k-1> \mid m < k \land k \in \mathbb{N}\}$とおくと,
$\BijectionMap{g}{\mathbb{N} \setminus \{m\}}{\mathbb{N}}$である.\\
$h:= f_{n+1} \setminus \{<n, m>\}$とおく.すると,
$\BijectionMap{h}{n}{\mathbb{N} \setminus \{m\}}$\\
よって、$\BijectionMap{g \circ h}{n}{\mathbb{N}}$\\
この写像の存在は帰納法の仮定に矛盾する.\\
\end{proof}

\begin{definition}[p184の定義7.1.3]
	\label{definition:合論的な自然数の定義}
	$X$を集合とする.\\
	1. $n$を自然数とする.
	全単射$n \rightarrow X$が存在するとき,
	$n$を$X$の濃度と呼び,
	$n=\operatorname{Card}(X)$と表す.\\
	2. 全単射$\mathbb{N} \rightarrow X$が存在するとき,
	$X$は加算無限集合であるといい,
	$\operatorname{Card}(X)=\aleph_0$と表す.\\
	$X$が有限集合であるか無限集合であるとき,
	$X$は加算であるという.
\end{definition}

\begin{proposition}[p185の命題7.1.5]
	部分集合$A \subset \mathbb{N}$に対し,
	関数$\map{C_A}{\mathbb{N}}{\mathbb{N}}$を帰納的に$C_A(0)=0$と,
  \begin{equation*}
    C_A(m + 1)=
    \begin{cases}
      C_A(m) + 1 & m \in A\text{のとき},\\
      C_A(m)  & m \not\in A \text{のとき}
    \end{cases}
  \end{equation*}
  で定める.\\
  1. $n$を自然数とし,
  $A \subseteq n$とする.
  このとき,
  $\operatorname{Card}(A)=C_A(n) \leq n$であり,
  $C_A$の制限は全単射$\map{C_A\restriction_{{A}}}{A}{\operatorname{Card}(A)}$を定める.
  また,
  $C_A(n)=n$ならば,
  $A=n$である.論理式は以下.\\
 $\forall n \in \mathbb{N}(A \subseteq n \rightarrow ((\operatorname{Card}(A)=C_A(n) \leq n ) \land (\BijectionMap{C_A\restriction_{{A}}}{A}{\operatorname{Card}(A)}
 ) \land (C_A(n) = n \rightarrow A=n)))$\\
 2. $A \subseteq n$をみたす自然数$n$は存在しないと仮定する.このとき,
 $\map {\seigen {C_A}{A}}{A}{\mathbb{N}}$は全単射である.論理式は以下.\\
 $\not \exists n (A \subseteq n) \rightarrow \BijectionMap{\seigen{C_A}{A}}{A}{\mathbb{N}}$
 
\end{proposition}

\begin{proof} $\varphi (n) :A \subseteq n \rightarrow ((\operatorname{Card}(A)=C_A(n) \leq n) \land (\BijectionMap{C_A\restriction_{{A}}}{A}{\operatorname{Card}(A)})
 \land (C_A(n) = n \rightarrow A=n))$とし、$n$に関する数学的帰納法で示す.\\
 
(base step) $\varphi (0) \Leftrightarrow A \subseteq 0 \rightarrow ((\operatorname{Card}(A)=C_A(0) \leq 0) \land (\BijectionMap{C_A\restriction_{{A}}}{A}{\operatorname{Card}(A)})
 \land (C_A(0) = n \rightarrow A=0)) $ \\
$A \subseteq 0$より,
$A=\emptyset$.
$\operatorname{Card}(A)=\operatorname{Card}(\emptyset)=\emptyset=C_A(0)$.\\
$C_A(0)\leq 0$も真より,
$\operatorname{Card}(A)=C_A(0) \leq 0$は真.\\
$A=\emptyset$より,
$\BijectionMap{C_A\restriction_{{A}}}{\emptyset}{\emptyset}$.今,
$\Card{A} = \emptyset$ より $\varphi_2(0)$は真.\\
$C_A(0)=0\rightarrow A=0$も成り立つ.\\

(induction step) $\forall n \in \mathbb{N} ( \varphi (n) \rightarrow \varphi (n+1)) \Leftrightarrow$ \\
$\forall n\in\mathbb{N} ((A \subseteq n \rightarrow ((\operatorname{Card}(A) = C_A(n) \leq n) \land (\BijectionMap{C_{A}\restriction_{{A}}}{A}{\operatorname{Card}(A)})
\land (C_A(n) = n \rightarrow A=n))) \\
 \rightarrow (A \subseteq n+1 \rightarrow ((\operatorname{Card}(A)=C_{A}(n+1) \leq n+1) \land (\BijectionMap{C_{A}\restriction_{{A}}}{A}{\operatorname{Card}(A)})
 \land (C_{A}(n+1) = n+1 \rightarrow A=n+1))))$\\
 任意に$n\in\mathbb{N}$をとり,
 前件を仮定して後件をしめす.\\
 $A \subseteq n+1$を仮定する.\\
 $A_n=A\cap n$とおくと,
 $A_n \subseteq A$であり,
 これと帰納法の仮定より$\varphi_1(n) \land \varphi_2(n) \land \varphi 3(n)$は真.\\
 また,
 $C_A \restriction {n+1} = C_{A_n} \restriction {n+1}$より,
 $C_A(n) = C_{A_n}(n)$.\\
 これと$C_A$の定義と$\varphi_1 (n)$より以下が成り立つ.\\
  \begin{equation*}
    C_A(n + 1)=
    \begin{cases}
     C_{A_n}(n) + 1 =  \operatorname{Card}(A_n) + 1 \leq n + 1& n \in A\text{のとき},\\
     C_{A_n}(n) = \operatorname{Card}(A_n) \leq n  & n \not\in A \text{のとき}
    \end{cases}
  \end{equation*}

 ($\varphi_1(n+1) \land \varphi_2(n+1)$)\\
 ($n \not\in A$) のとき\\
 $A=A_n$より,
 $\operatorname{Card}(A) = \operatorname{Card}(A_n).$\\
 よって,
 $C_A(n+1) = \operatorname{Card}(A_n) = \operatorname{Card}(A) \leq n \leq n+1.$\\
 $\varphi_2 (n)$と$A=A_n$より,
 $\varphi_2 (n+1)$は真.\\
 ($n \in A$) のとき\\
 $\operatorname{Card}(A) = \operatorname{Card}(A_n) + 1.$\\
 よって$C_A(n+1) = \operatorname{Card}(A_n) + 1 = \operatorname{Card}(A) \leq n+1$\\
 また,
 $A_n \subseteq n$と$C_A \restriction {n} = C_{A_n} \restriction {n}$より,
 $C_A \restriction {A_n} = C_{A_n} \restriction {A_n}$\\
 これと$\varphi_2(n)$より$\BijectionMap {C_A\restriction {A_n}}{A_n}{\Card{A_n}}$\\
 また,
 $C_A(n) = C_{A_n}(n) = \Card{A_n}$より,
 $\BijectionMap {\seigen {C_A}{A}}{A}{\Card{A}}$.\\
 
 $(\varphi_3(n+1))$\\
 $C_{A}(n+1)=n+1$と仮定する.\\
 このとき$n \in A$なので,
 $C_A(n+1) = C_{A_n}(n) + 1 = n + 1$より,
 $C_{A_n}(n) = n$.\\
 これと$\varphi_3(n)$より,
 $A_n = n$.\\
 $A = A_n \cup \{n\}$なので,
 $A = n+1$.\\
\end{proof}

\begin{proof}
 goal:= $\not \exists n (A \subseteq n) \rightarrow \InjectionMap{\seigen{C_A}{A}}{A}{\mathbb{N}} \land \SurjectionMap{\seigen{C_A}{A}}{A}{\mathbb{N}}$\\

 $\not \exists n (A \subseteq n)$が真とする.\\
 任意に$x_1, x_2 \in A$をとり,
 $\seigen {C_A}{A}(x_1) = \seigen{C_A}{A}(x_2)$とする.\\
 
 $(\InjectionMap{\seigen{C_A}{A}}{A}{\mathbb{N}}$)\\
 $m:= \operatorname{max}(x_1, x_2) + 1$とおき,
 $A_m := A \cap m$とおく.\\
 $x_1, x_2 < m$より,
 $x_1, x_2 \in m$.
 よって$x_1, x_2 \in A_m$\\
 また,
 $A_m \subseteq m$と1.より,
$\InjectionMap{\seigen {C_{A_m}}{A_m}}{A_m}{\Card{A_m}}$.\\
 これと$x_1, x_2 \in A_m$より,
 $\seigen{C_{A_m}}{A_m}(x_1) = \seigen{C_{A_m}}{A_m}(x_2) \rightarrow x_1 = x_2$.\\
 また,
 $\seigen {C_A}{m} = \seigen {C_{A_m}}{m}$と,
 $x_1, x_2 < m$より\\
 $\seigen {C_A}{A}(x_1) = \seigen {C_{A_m}}{A_m}(x_1), \seigen {C_A}{A}(x_1) = \seigen {C_{A_m}}{m}(x_2)$.\\
 よって$\seigen{C_{A_m}}{A_m}(x_1) = \seigen{C_{A_m}}{A_m}(x_2)$.
 なので,
 $x_1 = x_2$\\
 
 $(\SurjectionMap{\seigen{C_A}{A}}{A}{\mathbb{N}}$)\\

\end{proof}
